\begin{table}[htbp]
\centering
\small
\caption{Literature values used as targets and benchmarks.}
\label{tab:literature}
\begin{tabular}{lllll}
\hline
quantity & value & unit & source & used for \\
\hline
a (n-p triplet) & 5.411 & fm & Hackenburg, Phys. Rev. C 73, 044002 (2006) & Target of the deuteron tuning. Ref. [23] of RBEF 2023.\\
r0 (n-p triplet) & 1.744 & fm & Hackenburg, Phys. Rev. C 73, 044002 (2006) & Second target of the deuteron tuning.\\
E (deuteron) & -2.224 & MeV & Hackenburg, Phys. Rev. C 73, 044002 (2006) & Experimental benchmark for §4 — not an input.\\
a (4He dimer) & 90.4 & angstrom & Cencek et al., J. Chem. Phys. 136, 224303 (2012) & Target of the helium tuning. Ref. [22] of RBEF 2023.\\
r0 (4He dimer) & 8 & angstrom & Cencek et al., J. Chem. Phys. 136, 224303 (2012) & Second target of the helium tuning.\\
E (4He dimer) & -0.00162 & K & Cencek et al., J. Chem. Phys. 136, 224303 (2012) & Experimental benchmark for §4 — not an input.\\
a (n-n singlet) & -18.5 & fm & Macedo-Lima \& Madeira, Rev. Bras. Ensino Fis. 45, e20230079 (2023), Table 2 & Third tuning case: large NEGATIVE a, no bound state.\\
Published potential parameters & -- & - & Macedo-Lima \& Madeira (2023), Tables 3 and 4 & Initial guesses AND the validation target of §3.\\
v threshold, square well & 1.234 & - & Analytic: pi\textasciicircum{}2/8 = 1.2337 & Exact check of the first bound state in the §5 sweep. \\
\hline
\end{tabular}
\end{table}
